% BHU PYQ -- Professional Exam Template
\documentclass[11pt,a4paper]{article}

\usepackage[a4paper,top=2.5cm,bottom=2.5cm,left=2.8cm,right=2.8cm,headheight=14pt]{geometry}
\usepackage{amsmath,amssymb}
\usepackage[T1]{fontenc}
\usepackage[utf8]{inputenc}
\usepackage{lmodern}
\usepackage{microtype}
\usepackage{fancyhdr}
\usepackage{enumitem}
\usepackage{listings}
\usepackage{hyperref}
\usepackage{lastpage}
\usepackage{parskip}
\usepackage{tikz}

\hypersetup{colorlinks=true,urlcolor=black,linkcolor=black,
  pdftitle={BVCA-102: Fundamentals of Artificial Intelligence -- BHU PYQ}}

\pagestyle{fancy}
\fancyhf{}
\renewcommand{\headrulewidth}{0.4pt}
\renewcommand{\footrulewidth}{0.4pt}
\fancyhead[L]{\small   \textit{Banaras Hindu University}}
\fancyhead[R]{\small   \textit{BVCA-102: Fundamentals of Artificial Intelligence}}
\fancyfoot[C]{\small Page  \thepage\ of \pageref{LastPage}}

\lstset{
  basicstyle=\small tfamily,
  frame=single,
  breaklines=true,
  columns=flexible,
  keepspaces=true
}


\newlist{parts}{enumerate}{2}
\setlist[parts,1]{label=(\alph*),leftmargin=*,itemsep=4pt,topsep=3pt}
\setlist[parts,2]{label=(\roman*),leftmargin=*,itemsep=2pt,topsep=2pt}

\newcommand{\pts}[1]{\hfill   \textit{[#1]}}


\begin{document}

\begin{center}
  {\large\bfseries BANARAS HINDU UNIVERSITY}\\[2pt]
  {
\normalsize B.Voc. Semester I Examination, 2022-23}\\[6pt]
  \rule{0.8\linewidth}{0.5pt}\\[6pt]
  {\large\bfseries Paper: BVCA-102 --- Fundamentals of Artificial Intelligence}\\[4pt]
  {
\normalsize Time: Three Hours \hfill Maximum Marks: 70}
\end{center}

\rule{\linewidth}{0.4pt}

\smallskip



\noindent   \textit{Instructions: Answer All Sections. Candidates are required to write their answers in their own words as far as practicable.}

\bigskip

\smallskip

\bigskip


\noindent {\large\bfseries SECTION --- A}\\[4pt]
\rule{\linewidth}{0.4pt}\smallskip

Note: Long Answer type questions to be answered in about 750 words each. Attempt any two: \hfill   \textbf{15 $ \times$ 2 = 30 Marks}

\medskip


\noindent   \textbf{Question 1.}

\begin{parts}
  \item What is Artificial Intelligence? Discuss five applications of Artificial Intelligence.
  \item What are the challenges faced by Artificial Intelligence?
\end{parts}

\medskip


\noindent   \textbf{Question 2.}

\begin{parts}
  \item What are agents and environment? What is an Ideal Rational Agent?
  \item Discuss different types of agents in Artificial Intelligence.
\end{parts}

\medskip


\noindent   \textbf{Question 3.}

\begin{parts}
  \item What are Heuristic Search techniques? Discuss the Hill Climbing Algorithm.
  \item How does the Steepest Ascent Hill Climbing algorithm differ from the standard Hill Climbing algorithm?
\end{parts}

\medskip


\noindent   \textbf{Question 4.}

\begin{parts}
  \item Explain AO* algorithm with example.\pts{7}
  \item Discuss briefly A* algorithm. Find the optimal path and cost from Start node $S$ to Goal node $G$ for the given graph using A* algorithm. (The path cost is given on the edges of the graph and the table shows the values of heuristics cost).\pts{8}
  
  
\begin{center}
  
\begin{tikzpicture}[>=stealth,
    state/.style={draw, circle, minimum size=8mm, inner sep=0pt, font=\bfseries},
    weight/.style={draw, circle, fill=white, minimum size=5mm, inner sep=0pt, font=\small}
  ]
    
ode[state] (S) at (0,0) {S};
    
ode[state] (A) at (2,1.3) {A};
    
ode[state] (B) at (4.2,2.8) {B};
    
ode[state] (C) at (4.2,1.1) {C};
    
ode[state] (D) at (6.4,2.1) {D};
    
ode[state] (G) at (8.4,0) {G};
    
    \draw[->, thick] (S) -- node[weight] {1} (A);
    \draw[->, thick] (S) -- node[weight] {10} (G);
    \draw[->, thick] (A) -- node[weight] {2} (B);
    \draw[->, thick] (A) -- node[weight] {1} (C);
    \draw[->, thick] (B) -- node[weight] {5} (D);
    \draw[->, thick] (C) -- node[weight] {3} (D);
    \draw[->, thick] (C) -- node[weight] {4} (G);
    \draw[->, thick] (D) -- node[weight] {2} (G);
  \end{tikzpicture}
  \end{center}
  
  
\begin{center}
  
\begin{tabular}{|c|c|}
    \hline
   \textbf{State} &   \textbf{h(n)} \\
    \hline
    S & 5 \\
    \hline
    A & 3 \\
    \hline
    B & 4 \\
    \hline
    C & 2 \\
    \hline
    D & 6 \\
    \hline
    G & 0 \\
    \hline
  \end{tabular}
  \end{center}
\end{parts}

\bigskip


\noindent {\large\bfseries SECTION --- B}\\[4pt]
\rule{\linewidth}{0.4pt}\smallskip

Note: Short Answer type questions to be answered in
about 550 words each. (Any six) ($5  \times 6 = 30$)

\medskip


\noindent   \textbf{Question 5.} What are the steps for building an Expert System ?

\medskip


\noindent   \textbf{Question 6.} Discuss any five features of Environment.

\medskip


\noindent   \textbf{Question 7.} What do you understand by PEAS ? Design PEAS
information for Self-driving cars and Medical
Diagnostic System.

\medskip


\noindent   \textbf{Question 8.} What is Uninformed Search technique ? Discuss about
Depth First Search algorithm.

\medskip


\noindent   \textbf{Question 9.} Define Facts, Rules and Queries in terms of Prolog
Programming. Write a Prolog program to find
maximum of two numbers.

\medskip


\noindent   \textbf{Question 10.} What is Means End Analysis ? Explain with example.

\medskip


\noindent   \textbf{Question 11.} Differentiate between Forward Chaining and
Backward Chaining in the context of Expert Systems.

\medskip


\noindent   \textbf{Question 12.} What are Lists in prolog ? Define any three list
operations.

\bigskip


\noindent {\large\bfseries SECTION --- C}\\[4pt]
\rule{\linewidth}{0.4pt}\smallskip

\medskip


\noindent   \textbf{Question 13.} Select the correct option:\hfill ($1  \times 10 = 10$)

\begin{enumerate}[label=(\roman*),leftmargin=*,itemsep=12pt]
  \item What is state space?
  
\begin{enumerate}[label=(\alph*),itemsep=2pt]
    \item The whole problem
    \item Your Definition to a problem
    \item Problem you design
    \item Representing your problem with variables and parameters
  \end{enumerate}
  
  \item What is the space complexity of Depth-first search?
  
\begin{enumerate}[label=(\alph*),itemsep=2pt]
    \item $O(b)$
    \item $O(b^l)$
    \item $O(m)$
    \item $O(bm)$
  \end{enumerate}
  
  \item When is breadth-first search optimal?
  
\begin{enumerate}[label=(\alph*),itemsep=2pt]
    \item When there is less number of nodes
    \item When all step costs are equal
    \item When all step costs are unequal
    \item None of the mentioned
  \end{enumerate}
  
  \item The father of Artificial Intelligence is:
  
\begin{enumerate}[label=(\alph*),itemsep=2pt]
    \item Alan Turing
    \item John McCarthy
    \item Charles Babbage
    \item None of the above
  \end{enumerate}
  
  \item A technique that was developed to determine whether a machine could or could not demonstrate artificial intelligence is known as:
  
\begin{enumerate}[label=(\alph*),itemsep=2pt]
    \item Boolean Algebra
    \item Turing Test
    \item Logarithm
    \item Algorithm
  \end{enumerate}
  
  \item The component of an Expert system is:
  
\begin{enumerate}[label=(\alph*),itemsep=2pt]
    \item Knowledge Base
    \item Inference Engine
    \item User Interface
    \item All of the above
  \end{enumerate}
  
  \item An AI agent perceives and acts upon the environment using:
  
\begin{enumerate}[label=(\alph*),itemsep=2pt]
    \item Sensors
    \item Perceiver
    \item Actuators
    \item Both (a) and (c)
  \end{enumerate}
  
  \item The main function of a problem-solving agent is to:
  
\begin{enumerate}[label=(\alph*),itemsep=2pt]
    \item Solve the given problem and reach the goal
    \item Find out which sequence of actions will get it to the goal state
    \item Both (a) \& (b)
    \item None of the above
  \end{enumerate}
  
  \item Which of the following is not a variable in Prolog?
  
\begin{enumerate}[label=(\alph*),itemsep=2pt]
    \item Paul
    \item a\_123
    \item \_123
    \item Hello
  \end{enumerate}
  
  \item What would the output of the following Prolog query be?
  
\begin{verbatim}
  ?- f(a, b) = f(Y, X)
  \end{verbatim}
  
\begin{enumerate}[label=(\alph*),itemsep=2pt]
    \item X=b, Y=a.
    \item This expression will not compile
    \item X=a, Y=b.
    \item X=X, Y=Y.
  \end{enumerate}
\end{enumerate}

\vfill
\rule{\linewidth}{0.4pt}

\begin{center}\small
\href{https://bhu-exam.vercel.app/}{Click here to visit our website}\\[4pt]\href{https://me-aryan.vercel.app/}{Click here to visit our contributor}\\[4pt]\href{https://quashan-me.vercel.app/}{Click here to visit our contributor}
\end{center}
\end{document}